\section{Summary}
We show in section \ref{sec:HardFamilies} that standard message-passing graph neural networks (1-GNNs) are fundamentally unable to solve the clique-coloring principle, by leveraging known results in isomorphism testing. Explicitly, for every $k \ge 3$, we use graph products to construct an infinite family of regular graph pairs that are provably indistinguishable to any such GNN. Our theory holds also for higher order 2-GNNs. Complementing experiments show that both 1-GNNs and 2-GNNs could not achieve validation accuracy beyond random choice, while 3-GNNs achieved near perfect validation accuracy. 

In further experiments, we have shown under varying architectures and experimental specifications that CPP is invariably easy for standard GNNs over pairs of random graphs generated under the planted clique model (and analogous imposed k-coloring model). Our experiments show that distinguishing these graph pairs becomes non-trivial when graph pairs are drawn from distributions of matching average degree.